\scanpage{63}
\figure{\begin{tikzpicture}\draw(0,0)--(1.1,0);\end{tikzpicture}}

Avant de me lancer vers une exploitation des principes dégagés, il me faut encore développer le formalisme des ``points base à l'infini''. Le point de départ est le suivant. \deleted{\unsure{Et}} (NB je me borne pour l'instant à la dim 1, \deleted{bien que} \inserted{mais} il sera nécessaire de revenir par la suite sur les phénomènes analogues en dimension quelconque, relatifs à la donnée d'un diviseur à croisements normaux\ldots) Soit $X$ un schéma local régulier de dim 1 (i.e. le spectre d'un anneau \unsure{rég.} \inserted{$A$}), et $x$ son point fermé, $k=k(x)$,
\inserted{on suppose $k$ de car. nulle}
\deleted{$\unsure{U=}$}
\[
\tag{53} U=X\setminus\{x\}=\{\eta\}
\]
l'ouvert réduit au point générique, et
\[
\tag{54} T_x=\mathbf V(\mathfrak m_x/\mathfrak m_x^2)
=\check{\mathbf V}\bigl(\unsure{(\mathfrak m_x/\mathfrak m_x^2)|(\mathfrak m_x/\mathfrak m_x^2)^\vee}\bigr)
\]
le $k$-schéma ``espace tangent en $x$'', \inserted{$T_x^*=T_x\setminus\{0\}$,} je dis qu'on a alors un morphisme canonique
\[
\tag{55} \widehat\Pi_1(T_x^*)\longrightarrow\widehat\Pi_1(U)
\]
(où les $\widehat\Pi_1$ ici sont les $\Pi_1$ schématiques).

On est \unsure{ramené} \unsure{d'emblée} au cas où $A$ est \deleted{\unsure{\textup{à. A} est}} complet \deleted{\unsure{et alors}} (puisque $\widehat U\longrightarrow U$ donc $\widehat\Pi_1(\widehat U)\mathrel{-}\widehat\Pi_1(U)$)

Alors on sait qu'on a \unsure{!|:}
\mspage{62}
